\documentclass[11pt]{article}

% -------------------------------------------------
% Packages
% -------------------------------------------------
\usepackage{amsmath,amssymb,amsfonts}
\usepackage{geometry}
\usepackage{graphicx}
\usepackage{hyperref}
\usepackage{booktabs}

\geometry{margin=2.5cm}

\title{
	A Geometric Laplace Operator on Prime Quadruplets\\
	and Experimental Spectral Correlations with Riemann Zeros
}

\author{
	Thomas Hoffbauer
}

\date{\today}

\begin{document}
	
	\maketitle
	
	% -------------------------------------------------
	\begin{abstract}
		We study a geometrically defined operator on the discrete set of prime quadruplets $(p, p+2, p+6, p+8)$.
		
		By embedding these configurations into a logarithmic space, we define a weighted graph Laplacian that induces geometric similarity relations between quadruplets.
		
		Numerical experiments indicate a strong correlation between the spectrum of this operator and the imaginary parts of the nontrivial zeros of the Riemann zeta function.
		
		We emphasize that this observation is purely experimental. The purpose of this work is to define a precise operator framework and to investigate its spectral behavior in a reproducible setting.
	\end{abstract}
	
	% -------------------------------------------------
	\section{Introduction}
	
	The Hilbert--Pólya conjecture suggests that the nontrivial zeros of the Riemann zeta function may arise as the spectrum of a self-adjoint operator.
	
	A possible approach is to construct operators on arithmetically defined discrete structures and to study their spectral properties.
	
	In this work, we consider prime quadruplets
	\[
	(p, p+2, p+6, p+8),
	\]
	which form the densest admissible prime constellations.
	
	Rather than working with individual primes, we treat these quadruplets as structured objects and define a geometric operator on their configuration space.
	
	% -------------------------------------------------
	\section{Prime Quadruplets}
	
	For $N \ge 1$, let
	\[
	Q_N = \{q_1, \dots, q_N\}
	\]
	denote the first $N$ prime quadruplets, ordered by increasing initial prime.
	
	Each element
	\[
	q_k = (p_k, p_k+2, p_k+6, p_k+8)
	\]
	is defined by a prime $p_k$ such that all four numbers
	\[
	p_k,\; p_k+2,\; p_k+6,\; p_k+8
	\]
	are prime.
	
	% -------------------------------------------------
	\section{Geometric Embedding}
	
	We define an embedding
	\[
	x : Q_N \to \mathbb{R}^4,
	\quad
	x(q_k) = (\log p_k, \log(p_k+2), \log(p_k+6), \log(p_k+8)).
	\]
	
	This embedding maps each quadruplet to a point in Euclidean space.
	
	% -------------------------------------------------
	\section{Graph Laplacian}
	
	We define a weight matrix
	\[
	W_{ij} = \exp\left(-\frac{\|x_i - x_j\|^2}{\sigma^2}\right).
	\]
	
	The Gaussian kernel is chosen because:
	\begin{itemize}
		\item it yields a symmetric positive similarity matrix,
		\item it provides a controllable locality scale via $\sigma$,
		\item it is standard in spectral graph theory.
	\end{itemize}
	
	The degree matrix is
	\[
	D_{ii} = \sum_j W_{ij},
	\]
	
	and the Laplace operator is
	\[
	L = D - W.
	\]
	
	This operator is self-adjoint and positive semidefinite.
	
	% -------------------------------------------------
	\section{Extended Embedding}
	
	To incorporate internal structural information, we define an extended embedding:
	\[
	x^{(8)} =
	(\log p_1,\log p_2,\log p_3,\log p_4,
	\log\frac{p_2}{p_1},\log\frac{p_3}{p_2},
	\log\frac{p_4}{p_3},\log\frac{p_4}{p_1}).
	\]
	
	The additional coordinates encode dimensionless relative relationships within each quadruplet.
	
	% -------------------------------------------------
	\section{Spectral Comparison}
	
	Let
	\[
	0 = \lambda_1 \le \lambda_2 \le \cdots \le \lambda_N
	\]
	denote the eigenvalues of $L$.
	
	We consider the first $m$ nonzero eigenvalues:
	\[
	(\lambda_2, \dots, \lambda_{m+1}).
	\]
	
	Let
	\[
	(\gamma_1, \dots, \gamma_m)
	\]
	denote the imaginary parts of the first $m$ nontrivial zeros of $\zeta(s)$.
	
	Both sequences are affinely rescaled to mean zero and unit variance.
	
	We compute the Pearson correlation coefficient:
	\[
	\mathrm{corr}(\lambda_n, \gamma_n).
	\]
	
	% -------------------------------------------------
	\section{Numerical Results}
	
	Representative numerical experiments suggest correlations on the order of
	\[
	\mathrm{corr} \approx 0.9.
	\]
	
	A detailed parameter study and full numerical tables will be provided in the accompanying code repository.
	
	% -------------------------------------------------
	\section{Validation}
	
	We test robustness using:
	
	\begin{itemize}
		\item Shuffle test (random permutation of $Q_N$)
		\item Random baseline (random integer tuples)
		\item Parameter variation of $\sigma$
	\end{itemize}
	
	Preliminary experiments indicate that the observed correlation is not reproduced under these baseline modifications.
	
	% -------------------------------------------------
	\section{Discussion}
	
	The observed spectral similarity is notable but does not imply a structural connection between the operator and the zeta function.
	
	Key open questions include:
	\begin{itemize}
		\item whether the operator is canonical,
		\item whether a meaningful limit exists as $N \to \infty$,
		\item and whether connections to known analytic frameworks can be established.
	\end{itemize}
	
	The present work should therefore be viewed as exploratory.
	
	% -------------------------------------------------
	\section{Conclusion}
	
	We introduced a geometrically defined Laplace operator on prime quadruplets and investigated its spectral properties numerically.
	
	The results suggest that geometric approaches to structured prime configurations may provide a useful experimental framework for studying spectral phenomena related to the Riemann zeta function.
	
	% -------------------------------------------------
	\begin{thebibliography}{9}
		
		\bibitem{Montgomery}
		H. L. Montgomery,
		The pair correlation of zeros of the zeta function,
		Proc. Symp. Pure Math. 24 (1973).
		
		\bibitem{Odlyzko}
		A. Odlyzko,
		The $10^{20}$-th zero of the Riemann zeta function.
		
		\bibitem{Connes}
		A. Connes,
		Trace formula in noncommutative geometry,
		Selecta Mathematica (1999).
		
		\bibitem{Chung}
		F. Chung,
		Spectral Graph Theory,
		CBMS, AMS (1997).
		
		\bibitem{Belkin}
		M. Belkin and P. Niyogi,
		Laplacian eigenmaps,
		Neural Computation (2003).
		
	\end{thebibliography}
	
\end{document}